\section{Reproducibility and numerical safeguards}
\label{app:reproducibility}

\paragraph{Code and artifacts.}
The analysis code, run configurations, Lean library and manuscript sources are available at \url{https://github.com/ioannist/six-birds-neutrino}. The audit figures and the two numerical tables (\cref{tab:audit-summary,tab:archived-chains}) are generated directly from stored run bundles by \path{scripts/make_paper_figures.py}. Its output, \path{paper/figures/figure_sources.json}, records the plotted totals with their source bundles and the parameters of the analytic illustration. The two-constraint figure is evaluated afresh from the analytic formulas, using the same numerical routines as the toy-model runs and the parameters stated in its caption. Localization values and percentages quoted in the text come from the same bundles and from the multipole-range check listed below. The baseline table summarizes \cite{arxiv2601_16277}. The main bundles, all under \path{runs/}, are:
\begin{itemize}[leftmargin=*, itemsep=1pt]
  \item cosmological audits: \texttt{20261003\_\allowbreak{}math\_\allowbreak{}review\_\allowbreak{}cosmological\_\allowbreak{}spt\_\allowbreak{}desi}, \texttt{\dots\_\allowbreak{}spt\_\allowbreak{}desi\_\allowbreak{}staging}, \texttt{\dots\_\allowbreak{}cmb\_\allowbreak{}desi}, \texttt{\dots\_\allowbreak{}cmb\_\allowbreak{}desi\_\allowbreak{}warm}, \texttt{\dots\_\allowbreak{}cmb\_\allowbreak{}desi\_\allowbreak{}staging}, \texttt{\dots\_\allowbreak{}cmb\_\allowbreak{}desi\_\allowbreak{}tau\_\allowbreak{}control}, \texttt{\dots\_\allowbreak{}cmb\_\allowbreak{}desi\_\allowbreak{}accuracy} and \texttt{\dots\_\allowbreak{}cmb\_\allowbreak{}desi\_\allowbreak{}accuracy\_\allowbreak{}improved\_\allowbreak{}source};
  \item fixed-spectrum audits: \texttt{20261003\_\allowbreak{}math\_\allowbreak{}review\_\allowbreak{}profiled\_\allowbreak{}multistart}, \texttt{\dots\_\allowbreak{}profiled\_\allowbreak{}staging\_\allowbreak{}multistart} and \texttt{\dots\_\allowbreak{}verified\_\allowbreak{}native\_\allowbreak{}covariance};
  \item multipole-range totals: \texttt{20261003\_\allowbreak{}math\_\allowbreak{}review\_\allowbreak{}localization\_\allowbreak{}range\_\allowbreak{}check};
  \item archived chain diagnostics: \texttt{20261003\_\allowbreak{}math\_\allowbreak{}review\_\allowbreak{}chain\_\allowbreak{}diagnostics};
  \item two-constraint model (numerical checks of the routines used for \cref{fig:toy}): \texttt{20261003\_\allowbreak{}math\_\allowbreak{}review\_\allowbreak{}toy\_\allowbreak{}cancellation}.
\end{itemize}
Each cosmological bundle includes an independent re-evaluation of every endpoint with a freshly built model (\texttt{accounting\_\allowbreak{}verification.json}). A complete record of the checks behind this paper, including the multi-chain sampling campaign and its solver controls, is in \path{docs/findings/math_review.md}.

\paragraph{Capped nuisance sets in the fixed-spectrum audit.}
When D1 is the test lens, the re-fitted set is \texttt{Tcal\_\allowbreak{}ext150}, \texttt{Tcal\_\allowbreak{}rel90}, \texttt{Tcal\_\allowbreak{}rel220}, \texttt{Ecal\_\allowbreak{}ext150}, \texttt{Ecal\_\allowbreak{}rel90}, \texttt{Ecal\_\allowbreak{}rel220}, \texttt{TT\_\allowbreak{}CIBClustering\_Alpha}, \texttt{TT\_\allowbreak{}CIB\_150x150}, \texttt{TT\_\allowbreak{}CIB\_150x220}, \texttt{TT\_\allowbreak{}CIB\_220x220}, \texttt{TT\_\allowbreak{}GalCirrus\_Alpha} and \texttt{TT\_\allowbreak{}GalCirrus\_Amp}. When 2018 is the test lens, it is \texttt{TT\_\allowbreak{}tSZ\_CIB\_Corr\_Amp}, \texttt{TT\_\allowbreak{}tSZ\_CIB\_corr}, \texttt{Ecal150}, \texttt{Ecal220}, \texttt{Ecal90}, \texttt{TT\_\allowbreak{}CIBClustering\_Alpha}, \texttt{TT\_\allowbreak{}CIBClustering\_Amp}, \texttt{TT\_\allowbreak{}CIBClustering\_Beta}, \texttt{TT\_\allowbreak{}GalCirrus\_Alpha}, \texttt{TT\_\allowbreak{}GalCirrus\_Amp}, \texttt{TT\_\allowbreak{}GalCirrus\_Beta} and \texttt{TT\_\allowbreak{}Poisson\_150x150}. Both sets are capped at twelve parameters to keep the run time tractable. The shared parameters are the tSZ and kSZ amplitudes.

\paragraph{Numerical safeguards.}
The pipeline enforces the following properties, each covered by regression tests (245 Python tests in total):
\begin{itemize}[leftmargin=*, itemsep=1pt]
  \item \emph{Covariances} must be finite, symmetric and positive definite. Only round-off asymmetry is symmetrized, and singular matrices are rejected rather than pseudo-inverted. Solves, allocations and their sums that cannot be represented in floating point are refused rather than returned as infinities or NaNs.
  \item \emph{Quantiles} use exact integer or rational cumulative mass, so they do not depend on how samples are grouped into rows.
  \item \emph{Chains} are combined only when their likelihoods, priors, sampler temperature (which must be one), Boltzmann code version and numerical settings all match. Column names in chain files must be unique and unambiguous.
  \item \emph{Two-constraint model}: the combination of \cref{eq:toy_combined_gaussian}, the boundary probability of \cref{eq:sweep_probability} and the gated density remain accurate under extreme weights, near-cancelling means and far tails. They fall back to exact rational arithmetic, or to a bounded asymptotic tail expansion, when floating point would lose the answer.
  \item \emph{Toy template rewrite}: a covariance-weighted projection of a template out of a residual is computed in whitened coordinates, with exact accumulation. When the template is chosen from the residual itself, the improvement in fit is guaranteed by construction and is reported as descriptive only.
\end{itemize}
These safeguards are tested numerical properties, not floating-point error certificates.
